\begin{center}
\LARGE
Rapid Prototyping for the Construction of Higher Order
           Finite Element Methods on Sparse Grids 

\end{center}

\begin{center}
\Large
Hans-Joachim Bungartz


\end{center}

\noindent
Date:  July 18th (Thursday)\\
Time:  11:10-11:35\\

\begin{center}
\large
Abstract
\end{center}

For the development of new algorithmic concepts in the area of numerical analysis or, more
generally, scientific computing, modern tools for an efficient rapid prototyping like computer
algebra programs like Maple or Mathematica and shell-script-type interpreter languages like Perl
gain more and more in importance. This is mainly due to the fact that the period of time from an
idea to a first programmed version of the corresponding algorithm can be cut down significantly,
since all the implementation and declaration overhead coming along with standard (numerical)
programming languages can be avoided. As an example for that, we present a new unidirectional
approach for $d$-dimensional finite element methods of higher order on sparse grids that enables
us to deal with problems of an arbitrary number $d$ of dimensions and to use polynomial bases of
an arbitrary degree $p$ with the same storage requirements as and only a little bit more of
computational work than in the usual piecewise linear case. The code consists of two parts: a kind
of setup phase which is done in Maple and the iterative solver programmed in Perl. We report on
both numerical and implementational experiences and on interface problems resulting from the use
of tools originating from different ``worlds''. 

